% mixdom-wfst.tex — Labeled-MixDom Singlet HMM and WFST
% To be \input from tkf.tex (uses macros defined there)

\section{Labeled-MixDom Singlet HMM and WFST}
\label{sec:labeled-mixdom}

In this section we define state machines for the MixDom model that
foreground all latent state in the input/output alphabet, so they can
be used directly in beam-search MSA, progressive reconstruction, and
the variational ancestral-state framework of
Section~\ref{sec:varanc-presence-mixdom} without algebraic
distillation.
We call these the \emph{Labeled-MixDom Singlet HMM} and the
\emph{Labeled-MixDom WFST}, to distinguish them from the
Maraschino-distilled order-1 HMMs and WFSTs.

\begin{remark}[Within-fragment fragchar dynamics]
\label{rem:mixdom-fragchars}
Each fragment in MixDom carries an intra-fragment Markov chain on
fragchars: from a character with fragchar $\srcfrag$ in a fragment of
domain $\dom$, the next character (still in the same fragment) has
fragchar $\destfrag$ with probability $\ext^{(\dom)}_{\srcfrag\destfrag}$,
and the fragment terminates at the current character with probability
$\notext^{(\dom)}_\srcfrag = 1 - \sum_\destfrag \ext^{(\dom)}_{\srcfrag\destfrag}$.
This is consistent with the exploded-MixDom spec
(\S\ref{sec:mixdom-frag-trans} of exploded-mixdom.tex), where
$\ext^{(k)}_{fg}$ is the intra-fragment fragchar transition kernel.
The TKF92 \emph{fragment} concept is preserved ($g=1$ marks the last
character of a fragment).

\begin{itemize}[nosep]
\item \textbf{Within a fragment ($g=0$).} The next character is in
the same fragment with fragchar $\destfrag$ chosen via
$\ext^{(\dom)}_{\srcfrag\destfrag}$; the destination $\destfrag$ may
or may not equal the source $\srcfrag$ (intra-fragment Markov, not
self-loop only).

\item \textbf{Fragment boundary ($g=1$).} The current character is the
last in its fragment; this happens with the fragment-termination
weight $\notext^{(\dom)}_\srcfrag$. After fragment termination, the
TKF92 outer dynamics decide what happens next: with probability
$\kappa_\dom$ another fragment in the same domain begins (its first
character drawn from $\fragdist_{\dom, \cdot}$); with probability
$1-\kappa_\dom$ the domain ends ($e=1$).

\item \textbf{Domain termination ($e=1$).} The TKF91 outer dynamics
on domains (with parameter $\kappa_{\main}$) decide whether another
domain begins or the sequence ends.

\item \textbf{$\nfrag = 1$ special case.} Setting $N_\nfrag = 1$
(single fragchar) collapses $\ext^{(\dom)}$ to a scalar self-loop;
the resulting expressions reduce to a TKF92-extension scalar form
that admits a more compact algebra (used in the simpler closed forms
elsewhere in this paper).
\end{itemize}
\end{remark}

\subsection{Labeled Alphabet}

Each alphabet symbol $\anctok \in \alphabet$ is decorated with a label
tuple $(\class, \frag, g, \dom, e)$ consisting of the following latent
variables:
\begin{itemize}
\item $\class \in \nclass = \{1,\ldots,|\nclass|\}$: substitution site
class.
\item $\frag \in \nfrag = \{1,\ldots,|\nfrag|\}$: fragchar of the
current character (the per-character class label whose Markov chain
within a fragment is $\ext^{(\dom)}$).
\item $g \in \{0,1\}$: fragment-end indicator ($g=1$ if this is the
last character of the current fragment; $g=0$ otherwise).
\item $\dom \in \ndom = \{1,\ldots,|\ndom|\}$: domain type.
\item $e \in \{0,1\}$: domain-end indicator ($e=1$ if this is the
last fragment of the current domain).
\end{itemize}
We write $\anctok_{\class\frag g\dom e}$ for a labeled symbol and
define
$\mathcal{L} = \nclass \times \nfrag \times \{0,1\} \times \ndom \times \{0,1\}$,
so $|\mathcal{L}| = 4 |\nclass| |\nfrag| |\ndom|$. The labeled
alphabet has $|\alphabet| \cdot |\mathcal{L}|$ symbols.

For context tracking in order-1 machines we need only the
\emph{structural label} $\ell = (\frag, g, \dom, e)$, since the site
class $\class$ does not affect transition structure. The number of
distinct structural labels is $L = 4 |\nfrag| |\ndom|$.


\subsection{Labeled-MixDom Singlet HMM}
\label{sec:labeled-singlet}

The Labeled-MixDom Singlet HMM generates sequences from the MixDom stationary distribution
with all latent variables made explicit.
It is an order-1 HMM whose state records the structural label $\ell = (\frag, g, \dom, e)$
of the most recently emitted character.

\paragraph{States.}
The state space is $\{\sta\} \cup \{(\frag, g, \dom, e) : \frag \in \nfrag,\, g \in \{0,1\},\, \dom \in \ndom,\, e \in \{0,1\}\} \cup \{\fin\}$,
giving $L + 2$ states.

\paragraph{Emissions.}
In state $(\frag, g, \dom, e)$, the HMM emits symbol $\anctok_{\class\frag g\dom e}$
(with the $\frag, g, \dom, e$ components matching the state) with probability
\begin{equation}
\label{eq:singlet-emit}
\emprob(\anctok_{\class\frag g\dom e} \mid \frag, g, \dom, e)
= \classdist_{\frag\class}\, \eqm_{\class\anctok}
\end{equation}
summing to 1 over $(\class, \anctok)$.

\paragraph{Transitions.}
From the start state $\sta$:
\begin{equation}
\label{eq:singlet-start}
P(\sta \to (\frag, g, \dom, e)) = \domdist_\dom\, \kappa_\dom\, \fragdist_{\dom\frag}\,
\begin{cases}
1 - \notext^{(\dom)}_\frag & \text{if } g = 0 \\
\notext^{(\dom)}_\frag\, \kappa_\dom\, \fragdist_{\dom\frag'} & \text{if } g = 1,\, e = 0 \\
& \quad\text{(to state $(\frag', g', \dom, e')$; see below)} \\
\notext^{(\dom)}_\frag\,(1-\kappa_\dom)\, \kappa_\main\, \domdist_{\dom'}\, \kappa_{\dom'}\, \fragdist_{\dom'\frag'} & \text{if } g = 1,\, e = 1 \\
& \quad\text{(to state $(\frag', g', \dom', e')$; see below)}
\end{cases}
\end{equation}
However, since the singlet HMM is order-1 and we only track the \emph{destination} state,
it is cleaner to express the transitions directly.

The stationary distribution for the MixDom singlet process factors as
follows. At the domain level, a TKF91 process with parameters
$(\insrate_\main, \delrate_\main)$ generates
$N_{\rm dom} \sim \mathrm{Geom}(\kappa_\main)$ domains, each of type
$\dom \sim \catdist(\domdist_1,\ldots,\domdist_\ndom)$. Within domain
$\dom$, an irreducible Markov chain over fragchars with transition
matrix $\ext^{(\dom)} \in [0, 1]^{\nfrag \times \nfrag}$ and
termination weights
$\notext^{(\dom)}_\srcfrag = 1 - \sum_\destfrag \ext^{(\dom)}_{\srcfrag\destfrag}$
generates the sequence of per-character fragchars; the first character
in the domain is drawn from the chain's initial distribution
$\fragdist_{\dom\srcfrag}$. Within each character's fragchar
$\srcfrag$, the site class is $\class \sim \classdist_{\srcfrag\class}$
and the residue is $\anctok \sim \eqm_\class$.

\begin{center}
\renewcommand{\arraystretch}{1.4}
\begin{tabular}{lll}
\hline
Source $\ell$ & Dest $\ell'$ & Transition weight \\
\hline
$\sta$ & $(\frag', 0, \dom', e')$ &
  $\kappa_\main\, \domdist_{\dom'}\, \fragdist_{\dom'\frag'}$ \\
$\sta$ & $(\frag', 1, \dom', e')$ &
  $\kappa_\main\, \domdist_{\dom'}\, \fragdist_{\dom'\frag'}$ \\
$\sta$ & $\fin$ & $1 - \kappa_\main$ \\
\hline
\multicolumn{3}{l}{\emph{Mid-domain continuations:}} \\
$(\frag, 0, \dom, e)$ & $(\frag', g', \dom, e)$ &
  $\ext^{(\dom)}_{\frag\frag'}$ \\
$(\frag, 1, \dom, 0)$ & $(\frag', g', \dom, e')$ &
  $\notext^{(\dom)}_\frag \cdot \kappa_\dom \cdot \fragdist_{\dom\frag'}$ \\
$(\frag, 1, \dom, 1)$ & $(\frag', g', \dom', e')$ &
  $\notext^{(\dom)}_\frag \cdot (1-\kappa_\dom) \cdot \kappa_\main \cdot \domdist_{\dom'} \cdot \kappa_{\dom'} \cdot \fragdist_{\dom'\frag'}\,/\,(1-\zeta)$ \\
\hline
\multicolumn{3}{l}{\emph{Termination:}} \\
$(\frag, 0, \dom, e)$ & $\fin$ & $0$ \quad (cannot end mid-fragment) \\
$(\frag, 1, \dom, 0)$ & $\fin$ &
  $\notext^{(\dom)}_\frag \cdot (1-\kappa_\dom) \cdot (1-\kappa_\main)\,/\,(1-\zeta)$ \\
$(\frag, 1, \dom, 1)$ & $\fin$ &
  $\notext^{(\dom)}_\frag \cdot (1-\kappa_\dom) \cdot (1-\kappa_\main)\,/\,(1-\zeta)$ \\
\hline
\end{tabular}
\end{center}
The $g'$ and $e'$ in the destination labels range over all values
($g' \in \{0,1\}$, $e' \in \{0,1\}$); the destination's $g'$ is set by
the prior $P(g' = 1 | \frag', \dom') = \notext^{(\dom')}_{\frag'}$
(per-character fragment-termination probability under the
intra-fragment Markov), and likewise $e'$ by the per-fragment
domain-termination probability $1-\kappa_{\dom'}$. The weights above
are the joint contributions; conditional probabilities of next-state
indicators are recovered by multiplying by the appropriate prior
factor for $(g', e')$.

The $\zeta$ correction $\zeta = \kappa_\main \cdot
\emptyseg_0^{(\rm sing)}$ accounts for skipping over geometrically-many
empty domains before the next emitting character or sequence
termination, where $\emptyseg_0^{(\rm sing)} = \sum_\dom \domdist_\dom
(1-\kappa_\dom)$ is the singlet-process probability of an empty
domain (under per-domain TKF92 with parameter $\kappa_\dom$).
Verification of the marginalised $(\dom, \frag)$-only kernel
(Section~\ref{sec:varanc-presence-mixdom}, equation~\eqref{eq:omega})
against the existing MixDom Pair HMM implementation
\texttt{build\_nested\_trans} at $t=0$ confirms these weights to 5+
decimal places on a (n\_dom=2, n\_fr=2) test instance.

More precisely, the probability of ending the sequence from a domain boundary state
involves summing over all possible runs of empty domains before termination.
Let $\zeta = \kappa_\main \cdot \emptyseg_0^{(\rm sing)}$ be the probability of generating
an empty domain and continuing.
Then from state $(\frag, 1, \dom, 1)$ (domain boundary):
\begin{equation}
\label{eq:singlet-end-domain}
P(\text{to } \fin) = \frac{1 - \kappa_\main}{1 - \zeta}
\end{equation}
and from state $(\frag, 1, \dom, 0)$ (fragment boundary, mid-domain):
\begin{equation}
\label{eq:singlet-end-frag}
P(\text{to } \fin) = (1 - \kappa_\dom) \cdot \frac{1 - \kappa_\main}{1 - \zeta}
\end{equation}

Similarly, the transition from a fragment boundary to a new character must account for
possibly skipping empty domains. From $(\frag, 1, \dom, 1)$:
\begin{equation}
\label{eq:singlet-newdom}
P((\frag, 1, \dom, 1) \to (\frag', g', \dom', e'))
= \frac{\kappa_\main \cdot \domdist_{\dom'} \cdot \kappa_{\dom'} \cdot \fragdist_{\dom'\frag'}}{1 - \zeta}
\end{equation}
accounting for the possibility that some intermediate domains were empty
(the nonempty domain $\dom'$ is reached after a geometric number of empty-domain trials).

\paragraph{Normalization check.}
From state $(\frag, 1, \dom, 1)$, summing over all destinations including $\fin$:
\[
\sum_{\frag',g',\dom',e'} \frac{\kappa_\main \domdist_{\dom'} \kappa_{\dom'} \fragdist_{\dom'\frag'}}{1-\zeta}
+ \frac{1-\kappa_\main}{1-\zeta}
= \frac{\kappa_\main(1 - \emptyseg_0^{(\rm sing)}) + 1 - \kappa_\main}{1 - \zeta}
= \frac{1 - \zeta}{1 - \zeta} = 1.
\]
where we used $\sum_{\dom'} \domdist_{\dom'} \kappa_{\dom'} = 1 - \emptyseg_0^{(\rm sing)}$
and $\zeta = \kappa_\main \emptyseg_0^{(\rm sing)}$.
The sum over $g'$ and $e'$ is implicit in the fragment/domain generation that follows.


\subsection{Labeled-MixDom WFST}
\label{sec:labeled-wfst}

\begin{remark}[WFST tables use the matrix-kernel notation]
\label{rem:wfst-notation}
The transition tables below use the intra-fragment Markov kernel
$\ext^{(\dom)}_{\srcfrag\destfrag}$ (matrix indexed by source and
destination fragchar) and its per-row termination
$\notext^{(\dom)}_\frag = 1 - \sum_{\frag'} \ext^{(\dom)}_{\frag\frag'}$.
Where a within-fragment self-loop appears, the relevant entry is the
diagonal $\ext^{(\dom)}_{\frag\frag}$; where a new-fragchar-same-domain
move appears, it is an off-diagonal $\ext^{(\dom)}_{\frag\frag'}$ for
$\frag' \neq \frag$; the fragment-end weight is $\notext^{(\dom)}_\frag$.
\end{remark}

The Labeled-MixDom WFST represents the conditional distribution of a descendant
labeled sequence given an ancestral labeled sequence, separated by evolutionary time $\evoltime$.
When composed with the Labeled-MixDom Singlet HMM (Section~\ref{sec:labeled-singlet}),
it must reproduce the MixDom Pair HMM (Section~\ref{sec:nestpairhmm}).

\begin{remark}[Reduced kernel for variational inference]
\label{rem:reduced-wfst-pointer}
The full $(c, f, g, d, e)$ alphabet is what makes this WFST self-contained
(no algebraic distillation, exact for beam search and progressive reconstruction).
For the variational ancestral-state framework
(Appendix~\ref{sec:varanc-presence-mixdom}) the labelled WFST is marginalised
analytically over $(c, c', g, e, g', e')$ to a reduced per-character kernel
$\hat{T}_{ss'}((d,f),(d',f'); \branchlen, \theta)$ over just $(d, f)$, with
$3 N_\text{dom} N_\text{fr} + 2$ states. The marginalisation does \emph{not}
collapse to a single labelled-WFST entry: the per-character
labelled transition $(d, f) \to (d', f')$ admits up to three latent routes
(intra-fragment fragchar transition; new fragment, same domain; new domain
that may equal $d$), and the source's $(g, e)$ has a non-trivial posterior
over routes whenever $d' = d$. The reduced kernel is therefore a route-sum
$\hat{T}_{ss'} = \sum_r \omega^{(r)} \tilde{T}^{\text{lab},(r)}_{ss'}$
(eq.~\eqref{eq:T-hat}), which collapses to the cleaner form
$\omega \cdot \tilde{T}^{\text{lab}}$ only when the per-route labelled WFST
entries coincide --- a degenerate special case of the trivial
$\nfrag = 1, \ndom = 1$ instance.
Class marginalisation is trivial at the indel level since
the WFST indel block is class-independent. So the labelled WFST defined
here and the reduced kernel used in the variational appendix are the same
object viewed through different state-space lenses, with the variational
appendix providing the explicit route-decomposition.
\end{remark}

\paragraph{Design principles.}
The WFST has two kinds of states:
\begin{itemize}
\item \textbf{Emitting states} (``unready''): $\mat$, $\mathrm{I_F}$, $\mathrm{I_D}$, $\mathrm{D_F}$, $\mathrm{D_D}$.
  These consume an input character, produce an output character, or both, then make a mandatory
  null transition to a Wait state.
\item \textbf{Wait states} (``ready''): $\mathrm{W_M}$, $\mathrm{W_{D_F}}$, $\mathrm{W_{D_D}}$.
  These inspect the boundary indicators $(g, e)$ of the current context and choose the next
  emitting state based on the hierarchical boundary structure.
  Wait states also handle end-of-sequence.
\end{itemize}
In addition there are the non-emitting $\sta$ and $\fin$ states.

\paragraph{States.}
Each emitting or wait state carries a structural label context $\ell = (\frag, g, \dom, e)$.
The full state space is:
\[
\{\sta\} \;\cup\;
\{(X, \frag, g, \dom, e) : X \in \{\mat, \mathrm{I_F}, \mathrm{I_D}, \mathrm{D_F}, \mathrm{D_D},
  \mathrm{W_M}, \mathrm{W_{D_F}}, \mathrm{W_{D_D}}\}\}
\;\cup\; \{\fin\}
\]
Not all combinations occur (see constraints below), but the upper bound is $8L + 2$ states.

\paragraph{Context semantics.}
The context $\ell = (\frag, g, \dom, e)$ on each state records:
\begin{itemize}
\item In $\mat$ and $\mathrm{W_M}$: the label of the most recent matched character
  (which is the same for both input and output, since match preserves labels).
\item In $\mathrm{I_F}$, $\mathrm{I_D}$: the label of the most recent \emph{output} character.
\item In $\mathrm{D_F}$, $\mathrm{D_D}$, $\mathrm{W_{D_F}}$, $\mathrm{W_{D_D}}$: the label of the most recent \emph{input} character.
\end{itemize}
Thus only one context tuple is needed per state.


\subsubsection{\texorpdfstring{Emitting-State Transitions (Unready $\to$ Wait)}{Emitting-State Transitions (Unready -> Wait)}}
\label{sec:emit-to-wait}

Every emitting state makes a mandatory null transition (no input, no output) to its
corresponding wait state. These transitions carry the boundary-survival weights from the
nested TKF92/TKF91 structure.

\begin{center}
\renewcommand{\arraystretch}{1.5}
\begin{tabular}{lllll}
\hline
Source & Dest & Weight & Input & Output \\
\hline
$(\mat, \frag, g, \dom, e)$ & $(\mathrm{W_M}, \frag, g, \dom, e)$ &
  $w_{\mat}(g, e)$ & $\varepsilon$ & $\varepsilon$ \\
$(\mathrm{D_F}, \frag, g, \dom, e)$ & $(\mathrm{W_{D_F}}, \frag, g, \dom, e)$ &
  $w_{\mathrm{D_F}}(g, e)$ & $\varepsilon$ & $\varepsilon$ \\
$(\mathrm{D_D}, \frag, g, \dom, e)$ & $(\mathrm{W_{D_D}}, \frag, g, \dom, e)$ &
  $w_{\mathrm{D_D}}(g, e)$ & $\varepsilon$ & $\varepsilon$ \\
\hline
$(\mathrm{I_F}, \frag, g, \dom, e)$ & $(\mathrm{W_M}, \frag, g, \dom, e)$ &
  $w_{\mathrm{I_F}}(g, e)$ & $\varepsilon$ & $\varepsilon$ \\
$(\mathrm{I_D}, \frag, g, \dom, e)$ & $(\mathrm{W_M}, \frag, g, \dom, e)$ &
  $w_{\mathrm{I_D}}(g, e)$ & $\varepsilon$ & $\varepsilon$ \\
\hline
\end{tabular}
\end{center}

Wait: the insert states need more careful treatment.
An inserted fragment or domain is a complete sub-sequence generated by the descendant.
The insert states $\mathrm{I_F}$ and $\mathrm{I_D}$ handle character-level emissions
\emph{within} an inserted fragment, and upon fragment/domain completion, control returns
to the wait state that initiated the insertion.
We therefore need to track whether we are inserting at the fragment level or domain level.

Let us reconsider the state structure more carefully.

\subsubsection{Revised State Structure}
\label{sec:revised-states}

In the MixDom Pair HMM, the five state types $\mat\mat$, $\mat\ins$, $\mat\del$, $\ins\ins$, $\del\del$
at each $(\dom, \frag)$ position represent:
\begin{itemize}
\item $\mat\mat_{\dom\frag}$: ancestral and descendant both have a character (match/substitution)
\item $\mat\ins_{\dom\frag}$: descendant insertion within domain $\dom$, fragment $\frag$
\item $\mat\del_{\dom\frag}$: ancestral deletion within domain $\dom$, fragment $\frag$
\item $\ins\ins_{\dom\frag}$: insertion of an entire domain (both ancestor and descendant insert)
\item $\del\del_{\dom\frag}$: deletion of an entire domain (both ancestor and descendant delete)
\end{itemize}

The top-level states $\mat, \ins, \del$ refer to the \emph{domain-level} TKF91 process,
while the nested states $\mat, \ins, \del$ refer to the \emph{fragment-level} TKF92 process
within a domain.

For the WFST, we separate the domain-level and fragment-level indel processes:

\begin{center}
\renewcommand{\arraystretch}{1.3}
\begin{tabular}{lll}
\hline
WFST State Type & Input & Output \\
\hline
$\mat$ (Match) & $\anctok_{\class\frag g\dom e}$ & $\destok_{\class\frag g\dom e}$ \\
$\mathrm{I_F}$ (Insert Fragment char) & $\varepsilon$ & $\destok_{\class\frag g\dom e}$ \\
$\mathrm{I_D}$ (Insert Domain char) & $\varepsilon$ & $\destok_{\class\frag g\dom e}$ \\
$\mathrm{D_F}$ (Delete Fragment char) & $\anctok_{\class\frag g\dom e}$ & $\varepsilon$ \\
$\mathrm{D_D}$ (Delete Domain char) & $\anctok_{\class\frag g\dom e}$ & $\varepsilon$ \\
$\mathrm{W_M}$ (Wait after Match) & $\varepsilon$ & $\varepsilon$ \\
$\mathrm{W_{D_F}}$ (Wait after Delete-Fragment) & $\varepsilon$ & $\varepsilon$ \\
$\mathrm{W_{D_D}}$ (Wait after Delete-Domain) & $\varepsilon$ & $\varepsilon$ \\
\hline
\end{tabular}
\end{center}

\paragraph{Key constraint: label preservation.}
In a Match state, the WFST does not change the structural label: the input and output
labels $(\frag, g, \dom, e)$ must be identical (though $\class$ is also preserved and
the character $\anctok \to \destok$ may change via substitution).
Insertions create \emph{new} characters with new labels; deletions consume characters
without producing output.


\subsubsection{Emitting to Wait Transitions}
\label{sec:emit-wait-transitions}

After emitting (or consuming), each emitting state transitions to its wait state.
These null transitions carry weights that account for the fragment-extension
and domain-continuation structure.

In the MixDom model, within a domain of type $\dom$, each fragment evolves
under the intra-fragment Markov kernel $\ext^{(\dom)}_{\srcfrag\destfrag}$.
Within the fragment ($g=0$), the character simply continues.
At a fragment boundary ($g=1$), the TKF92 process within the domain decides
whether to start a new fragment or end the domain.
At a domain boundary ($g=1, e=1$), the TKF91 domain-level process decides
whether to start a new domain or end the sequence.

The weights on emitting $\to$ wait transitions are:

\begin{center}
\renewcommand{\arraystretch}{1.5}
\begin{tabular}{lll}
\hline
Transition & Condition & Weight \\
\hline
$\mat \to \mathrm{W_M}$ & $g=0$ & $1$ \\
$\mat \to \mathrm{W_M}$ & $g=1, e=0$ &
  $(1-\beta_\dom)$ \\
$\mat \to \mathrm{W_M}$ & $g=1, e=1$ &
  $(1-\beta_\dom)(1-\beta_\main)$ \\[4pt]
$\mathrm{D_F} \to \mathrm{W_{D_F}}$ & $g=0$ & $1$ \\
$\mathrm{D_F} \to \mathrm{W_{D_F}}$ & $g=1, e=0$ &
  $(1-\gamma_\dom)$ \\
$\mathrm{D_F} \to \mathrm{W_{D_F}}$ & $g=1, e=1$ &
  $(1-\gamma_\dom)(1-\beta_\main)$ \\[4pt]
$\mathrm{D_D} \to \mathrm{W_{D_D}}$ & $g=0$ & $1$ \\
$\mathrm{D_D} \to \mathrm{W_{D_D}}$ & $g=1, e=0$ & $1$ \\
$\mathrm{D_D} \to \mathrm{W_{D_D}}$ & $g=1, e=1$ &
  $(1-\gamma_\main)$ \\
\hline
\end{tabular}
\end{center}
where $\beta_\dom = \beta(\insrate_\dom, \delrate_\dom, \evoltime)$,
$\gamma_\dom = \gamma(\insrate_\dom, \delrate_\dom, \evoltime)$,
$\beta_\main = \beta(\insrate_\main, \delrate_\main, \evoltime)$,
$\gamma_\main = \gamma(\insrate_\main, \delrate_\main, \evoltime)$.

The rationale: at $g=0$ we are mid-fragment, so no boundary weight is needed.
At $g=1$ (fragment boundary), the TKF92 boundary weight $(1-\beta)$ or $(1-\gamma)$ applies.
At $g=1, e=1$ (domain boundary), both the fragment boundary and the domain boundary
weights apply.
For $\mathrm{D_D}$, the domain is being deleted as a unit; within the deleted domain,
fragment structure is irrelevant (all fragments are consumed), so the fragment-level
weights are unity and only the domain-level weight $(1-\gamma_\main)$ applies at domain end.

\paragraph{Insert states.}
Insert states ($\mathrm{I_F}$, $\mathrm{I_D}$) represent characters being inserted
in the descendant.  An inserted fragment is a self-contained TKF92 fragment;
an inserted domain is a self-contained TKF91 domain.

After emitting an inserted character, the insert state loops back to itself
(fragment extension) or transitions to a wait state (fragment/domain termination).
The fragment extension self-loop:
\begin{center}
\renewcommand{\arraystretch}{1.3}
\begin{tabular}{llll}
\hline
Transition & Weight & Input & Output \\
\hline
$(\mathrm{I_F}, \srcfrag, g, \dom, e) \to (\mathrm{I_F}, \destfrag, g', \dom, e)$ &
  $\ext^{(\dom)}_{\srcfrag\destfrag}$ & $\varepsilon$ & $\destok_{\class\destfrag g'\dom e}$ \\
$(\mathrm{I_D}, \srcfrag, g, \dom, e) \to (\mathrm{I_D}, \destfrag, g', \dom, e)$ &
  $\ext^{(\dom)}_{\srcfrag\destfrag}$ & $\varepsilon$ & $\destok_{\class\destfrag g'\dom e}$ \\
\hline
\end{tabular}
\end{center}

On fragment termination within an inserted domain ($\mathrm{I_D}$), a new fragment
may begin (with TKF92 parameters for the inserted domain):
\begin{center}
\renewcommand{\arraystretch}{1.3}
\begin{tabular}{llll}
\hline
Transition & Weight & Condition & I/O \\
\hline
$(\mathrm{I_D}, \frag, 1, \dom, 0) \to (\mathrm{I_D}, \frag', g', \dom, e')$ &
  $\notext^{(\dom)}_\frag \cdot \kappa_\dom \cdot \fragdist_{\dom\frag'}$ &
  new frag in inserted domain &
  $\varepsilon / \destok$ \\
$(\mathrm{I_D}, \frag, 1, \dom, e) \to (\mathrm{W_M}, \ldots)$ &
  $\notext^{(\dom)}_\frag(1-\kappa_\dom)$ &
  $e$ set appropriately &
  $\varepsilon / \varepsilon$ \\
\hline
\end{tabular}
\end{center}

For $\mathrm{I_F}$ (inserted fragment within an existing domain), fragment termination
returns control to $\mathrm{W_M}$:
\begin{center}
\renewcommand{\arraystretch}{1.3}
\begin{tabular}{lll}
\hline
Transition & Weight & Condition \\
\hline
$(\mathrm{I_F}, \frag, 1, \dom, e) \to (\mathrm{W_M}, \frag, 1, \dom, e)$ &
  $\notext^{(\dom)}_\frag$ & fragment ends \\
\hline
\end{tabular}
\end{center}


\subsubsection{\texorpdfstring{Wait-State Transitions (Ready $\to$ Emitting)}{Wait-State Transitions (Ready -> Emitting)}}
\label{sec:wait-transitions}

Wait states inspect the boundary indicators and decide the next action.
The transitions depend on whether we are mid-fragment ($g=0$),
at a fragment boundary ($g=1, e=0$), or at a domain boundary ($g=1, e=1$).

\paragraph{Case 1: Mid-fragment ($g = 0$).}
Within a fragment, only continuation of the current fragment is possible.
No new fragments or domains can start.

\begin{center}
\renewcommand{\arraystretch}{1.5}
\begin{tabular}{lllll}
\hline
Source & Dest & Weight & Input & Output \\
\hline
$(\mathrm{W_M}, \frag, 0, \dom, e)$ &
  $(\mat, \frag, g', \dom, e)$ &
  $\alpha_\dom$ &
  $\anctok_{\class\frag g'\dom e}$ &
  $\destok_{\class\frag g'\dom e}$ \\
$(\mathrm{W_M}, \frag, 0, \dom, e)$ &
  $(\mathrm{D_F}, \frag, g', \dom, e)$ &
  $(1-\alpha_\dom)$ &
  $\anctok_{\class\frag g'\dom e}$ &
  $\varepsilon$ \\
\hline
$(\mathrm{W_{D_F}}, \frag, 0, \dom, e)$ &
  $(\mathrm{D_F}, \frag, g', \dom, e)$ &
  $1$ &
  $\anctok_{\class\frag g'\dom e}$ &
  $\varepsilon$ \\
\hline
$(\mathrm{W_{D_D}}, \frag, 0, \dom, e)$ &
  $(\mathrm{D_D}, \frag, g', \dom, e)$ &
  $1$ &
  $\anctok_{\class\frag g'\dom e}$ &
  $\varepsilon$ \\
\hline
\end{tabular}
\end{center}
Here $g'$ can be 0 or 1 (determined by the input character's label),
$\alpha_\dom = \alpha(\insrate_\dom, \delrate_\dom, \evoltime)$,
and the fragment type $\frag$ and domain indicators $(\dom, e)$ are unchanged.
The emission weight for $\mat$ is $\exp(\revsub_\class \evoltime)_{\anctok\destok}$;
the emission weight for $\mathrm{D_F}$ is 1 (input consumed, no output).

\paragraph{Case 2: Fragment boundary, mid-domain ($g = 1, e = 0$).}
At a fragment boundary within a domain, the TKF92 process within the domain
decides: start a new fragment (match or delete), or insert a new fragment.

\begin{center}
\renewcommand{\arraystretch}{1.5}
\begin{tabular}{lllll}
\hline
Source & Dest & Weight & Input & Output \\
\hline
$(\mathrm{W_M}, \frag, 1, \dom, 0)$ &
  $(\mat, \frag', g', \dom, e')$ &
  $\alpha_\dom \cdot \fragdist_{\dom\frag'}$ &
  $\anctok_{\class'\frag' g'\dom e'}$ &
  $\destok_{\class'\frag' g'\dom e'}$ \\
$(\mathrm{W_M}, \frag, 1, \dom, 0)$ &
  $(\mathrm{D_F}, \frag', g', \dom, e')$ &
  $(1-\alpha_\dom) \cdot \fragdist_{\dom\frag'}$ &
  $\anctok_{\class'\frag' g'\dom e'}$ &
  $\varepsilon$ \\
$(\mathrm{W_M}, \frag, 1, \dom, 0)$ &
  $(\mathrm{I_F}, \frag', g', \dom, 0)$ &
  $\beta_\dom \cdot \fragdist_{\dom\frag'}$ &
  $\varepsilon$ &
  $\destok_{\class'\frag' g'\dom 0}$ \\
$(\mathrm{W_M}, \frag, 1, \dom, 0)$ &
  $\fin$ &
  $(1-\kappa_\dom) \cdot (1-\kappa_\main)$ &
  $\varepsilon$ &
  $\varepsilon$ \\
\hline
$(\mathrm{W_{D_F}}, \frag, 1, \dom, 0)$ &
  $(\mat, \frag', g', \dom, e')$ &
  $\alpha_\dom \cdot \fragdist_{\dom\frag'}$ &
  $\anctok_{\class'\frag' g'\dom e'}$ &
  $\destok_{\class'\frag' g'\dom e'}$ \\
$(\mathrm{W_{D_F}}, \frag, 1, \dom, 0)$ &
  $(\mathrm{D_F}, \frag', g', \dom, e')$ &
  $(1-\alpha_\dom) \cdot \fragdist_{\dom\frag'}$ &
  $\anctok_{\class'\frag' g'\dom e'}$ &
  $\varepsilon$ \\
\hline
$(\mathrm{W_{D_D}}, \frag, 1, \dom, 0)$ &
  $(\mathrm{D_D}, \frag', g', \dom, e')$ &
  $\fragdist_{\dom\frag'}$ &
  $\anctok_{\class'\frag' g'\dom e'}$ &
  $\varepsilon$ \\
\hline
\end{tabular}
\end{center}

Note: in this case, a new fragment type $\frag'$ is drawn from $\fragdist_{\dom\frag'}$,
and $e'$ is determined by the input character's label.
The domain type $\dom$ is unchanged.

\paragraph{Normalization of $\mathrm{W_M}$ at fragment boundary ($g=1, e=0$).}
The outgoing weights from $(\mathrm{W_M}, \frag, 1, \dom, 0)$ must sum to 1 when
we include all possible input/output characters:
\begin{itemize}
\item With an input character present (ancestral fragment continues):
  $\alpha_\dom + (1 - \alpha_\dom) = 1$, weighted by $\fragdist_{\dom\frag'}$.
  But the input character determines $\frag'$, so the $\fragdist$ weight is a prior
  for the singlet composition, not a transition weight in the WFST.
\item With no input character:
  insertion weight $\beta_\dom$ or end weight.
\end{itemize}

Actually, for a WFST, the normalization is more subtle: the transducer weights need not
sum to 1 at every state, because the transducer represents a conditional distribution.
However, when composed with the singlet HMM, the resulting Pair HMM transitions must
be properly normalized.

\paragraph{Global vs.\ local normalization.}
The Labeled-MixDom WFST is constructed by dividing the row-stochastic Pair HMM
$\transnest$ (Section~\ref{sec:nestpairhmm})
by the row-stochastic Singlet HMM (Section~\ref{sec:labeled-singlet}).
By construction, $\text{Singlet} \circ \text{WFST} = \text{Pair HMM}$, so the
WFST is conditionally normalized in the global sense:
$\sum_{\text{output sequences } y} P(y \mid x, \theta, \evoltime) = 1$
for any ancestor sequence $x$.
However, the per-state outgoing weights from $\mat,\mathrm{I_F},\mathrm{I_D},
\mathrm{D_F},\mathrm{D_D}$ do \emph{not} sum to~1 over a fixed input symbol, even
when the Singlet emission factor for the destination label is included.
This is the same state-folding artifact discussed for the TKF92 WFST in
Appendix~\ref{sec:wfst-row-sums}: the Bernoulli-$\ext$ extension-vs-exit decision
(here, the $\ext^{(\dom)}_{\srcfrag\destfrag}$ vs.\ $\notext^{(\dom)}_\srcfrag$
event), and the $\kappa_\dom$ vs.\ $1{-}\kappa_\dom$ domain-continuation event,
have been compiled into the same $\mat\to\cdot$ and $\mathrm{I_\cdot}\to\cdot$
edges that already carry the destination-singlet factors.
Splitting each emitting state into a ``just-arrived'' and a ``decision'' state
would restore local stochasticity at the cost of doubling (or tripling) the state
graph. The compact form here trades local stochasticity for a smaller machine,
exactly as in TKF92.

Let us instead directly specify the transition weights that, when composed with the
Singlet HMM, reproduce the Pair HMM transition matrix $\transnest$ from
Section~\ref{sec:nestpairhmm}.

\subsubsection{Complete WFST Transition Table}
\label{sec:wfst-complete}

We now give the complete transition table.
For readability, we abbreviate the state labels and split by source wait-state type and
boundary case.

Let the following shorthand apply throughout:
\begin{align*}
\alpha_\dom &= \alpha(\insrate_\dom, \delrate_\dom, \evoltime), &
\beta_\dom &= \beta(\insrate_\dom, \delrate_\dom, \evoltime), &
\gamma_\dom &= \gamma(\insrate_\dom, \delrate_\dom, \evoltime) \\
\alpha_\main &= \alpha(\insrate_\main, \delrate_\main, \evoltime), &
\beta_\main &= \beta(\insrate_\main, \delrate_\main, \evoltime), &
\gamma_\main &= \gamma(\insrate_\main, \delrate_\main, \evoltime) \\
\kappa_\dom &= \insrate_\dom / \delrate_\dom, &
\kappa_\main &= \insrate_\main / \delrate_\main
\end{align*}
and recall $\nonemptytrans$ is the effective $5 \times 5$ domain-level transition matrix
with null states eliminated (Section~\ref{sec:nestpairhmm}).

\paragraph{Start transitions.}
From $\sta$, the WFST enters its first emitting state.
The weights mirror the first row of $\transnest$:

\begin{center}
\renewcommand{\arraystretch}{1.4}
\begin{tabular}{llll}
\hline
Dest & Weight & Input & Output \\
\hline
$(\mat, \frag', g', \dom', e')$ &
  $\nonemptytrans_{\sta\mat} \cdot \domdist_{\dom'} \cdot \tkftrans^{(\dom')}_{\sta\ystate} \cdot \fragdist_{\dom'\frag'}$ &
  $\anctok_{\class'\frag' g'\dom' e'}$ &
  $\destok_{\class'\frag' g'\dom' e'}$ \\
$(\mathrm{I_D}, \frag', g', \dom', e')$ &
  $\nonemptytrans_{\sta\ins} \cdot \domdist_{\dom'} \cdot \kappa_{\dom'} \cdot \fragdist_{\dom'\frag'}$ &
  $\varepsilon$ &
  $\destok_{\class'\frag' g'\dom' e'}$ \\
$(\mathrm{D_D}, \frag', g', \dom', e')$ &
  $\nonemptytrans_{\sta\del} \cdot \domdist_{\dom'} \cdot \kappa_{\dom'} \cdot \fragdist_{\dom'\frag'}$ &
  $\anctok_{\class'\frag' g'\dom' e'}$ &
  $\varepsilon$ \\
$\fin$ &
  $\nonemptytrans_{\sta\fin}$ &
  $\varepsilon$ & $\varepsilon$ \\
\hline
\end{tabular}
\end{center}
where $\ystate$ is determined by the nested state type of the destination
($\mat$ for Match, $\ins$ for Insert-Fragment, $\del$ for Delete-Fragment within the domain).
For Match destinations that enter at the first character of a domain,
$\tkftrans^{(\dom')}_{\sta\ystate}$ is the TKF92 Pair HMM transition from $\sta$
into the appropriate nested state.

\paragraph{Emitting to Wait transitions.}
After each emitting state, a null transition to the corresponding wait state occurs.
The weight depends on the boundary indicators:

\begin{center}
\renewcommand{\arraystretch}{1.4}
\begin{tabular}{lll}
\hline
Transition & Condition on $(g, e)$ & Weight \\
\hline
\multirow{3}{*}{$(\mat, \ell) \to (\mathrm{W_M}, \ell)$}
  & $g=0$ & $1$ \\
  & $g=1,\, e=0$ & $\notext^{(\dom)}_\frag \cdot 1$ \\
  & $g=1,\, e=1$ & $\notext^{(\dom)}_\frag \cdot 1$ \\[2pt]
\multirow{3}{*}{$(\mathrm{D_F}, \ell) \to (\mathrm{W_{D_F}}, \ell)$}
  & $g=0$ & $1$ \\
  & $g=1,\, e=0$ & $\notext^{(\dom)}_\frag$ \\
  & $g=1,\, e=1$ & $\notext^{(\dom)}_\frag$ \\[2pt]
\multirow{3}{*}{$(\mathrm{D_D}, \ell) \to (\mathrm{W_{D_D}}, \ell)$}
  & $g=0$ & $1$ \\
  & $g=1,\, e=0$ & $\notext^{(\dom)}_\frag$ \\
  & $g=1,\, e=1$ & $\notext^{(\dom)}_\frag$ \\
\hline
\end{tabular}
\end{center}

Wait---the fragment extension must also be handled.
At $g=0$, the character is mid-fragment; the next character in the same
fragment follows with certainty (the matrix entry $\ext^{(\dom)}_{\srcfrag\destfrag}$
is already accounted for in the singlet HMM emission of the next labelled
character).
At $g=1$, the fragment has ended, so the per-source termination weight
$\notext^{(\dom)}_\frag$ has already been ``spent'' by the fact that $g=1$
was observed.
Since the boundary indicators are part of the \emph{label on the character},
which is determined by the input (for ancestral) or output (for descendant),
the fragment extension probability is absorbed into the singlet HMM, not the WFST.

Therefore all emitting-to-wait transitions have \textbf{unit weight}:
\begin{equation}
\label{eq:emit-to-wait}
w(X \to W_X) = 1 \quad \text{for all } X \in \{\mat, \mathrm{D_F}, \mathrm{D_D}\}
\end{equation}
The fragment and domain boundary structure is encoded in the
\emph{wait-state outgoing transitions}, which condition on $(g, e)$.


\subsubsection{Wait-State Outgoing Transitions}
\label{sec:wait-outgoing}

The wait states make all structural decisions.
We organize by source state type and boundary case.

\paragraph{$\mathrm{W_M}$ (Wait after Match).}

\bigskip
\noindent\textbf{Case $g=0$ (mid-fragment):}
Continue the current fragment. The next input character must have the same $(\frag, \dom, e)$.

\begin{center}
\renewcommand{\arraystretch}{1.4}
\begin{tabular}{lllll}
\hline
Dest & Weight & Input & Output & Notes \\
\hline
$(\mat, \frag, g', \dom, e)$ &
  $\alpha_\dom$ &
  $\anctok_{\class'\frag g'\dom e}$ &
  $\destok_{\class'\frag g'\dom e}$ &
  match continues \\
$(\mathrm{D_F}, \frag, g', \dom, e)$ &
  $(1-\alpha_\dom)$ &
  $\anctok_{\class'\frag g'\dom e}$ &
  $\varepsilon$ &
  fragment-level deletion \\
\hline
\end{tabular}
\end{center}

\noindent\textbf{Case $g=1, e=0$ (fragment boundary, mid-domain):}
Fragment ended, new fragment within same domain, or insert/delete fragment,
or end domain and transition at domain level.

\begin{center}
\renewcommand{\arraystretch}{1.4}
\begin{tabular}{lllll}
\hline
Dest & Weight & I & O & Notes \\
\hline
$(\mat, \frag', g', \dom, e')$ &
  $\tkftrans^{(\dom)}_{\mat\mat} \cdot \fragdist_{\dom\frag'}$ &
  $\anctok_{\ldots}$ &
  $\destok_{\ldots}$ &
  new frag, match \\
$(\mathrm{D_F}, \frag', g', \dom, e')$ &
  $\tkftrans^{(\dom)}_{\mat\del} \cdot \fragdist_{\dom\frag'}$ &
  $\anctok_{\ldots}$ &
  $\varepsilon$ &
  new frag, delete \\
$(\mathrm{I_F}, \frag', g', \dom, 0)$ &
  $\tkftrans^{(\dom)}_{\mat\ins} \cdot \fragdist_{\dom\frag'}$ &
  $\varepsilon$ &
  $\destok_{\ldots}$ &
  insert frag \\
$\fin_{\rm domain}$ &
  $\tkftrans^{(\dom)}_{\mat\fin}$ &
  &
  &
  domain ends (see below) \\
\hline
\end{tabular}
\end{center}

When the domain ends, $\tkftrans^{(\dom)}_{\mat\fin} = (1-\beta_\dom)(1-\kappa_\dom)$, the
domain-level process takes over.
This is where the $\nonemptytrans$ matrix from the null-eliminated domain-level
Pair HMM applies.

Rather than introducing an intermediate ``domain end'' state, we can fold the
domain-level transition into the fragment-boundary transitions.
From $(\mathrm{W_M}, \frag, 1, \dom, 0)$, with the domain ending:
the TKF92 weight is $\tkftrans^{(\dom)}_{\mat\fin}$, and then the domain-level
$\nonemptytrans_{\mat\cdot}$ applies to reach the next domain.

However, this conflates two levels of the hierarchy.
To keep the WFST clean and composable, we should handle this through the
domain-end indicator $e$: when $e=0$, we are mid-domain and only fragment-level
transitions apply.  The domain-end case $e=1$ is reached when the last fragment
of a domain has ended.

Actually, let us reconsider.
The labels $(g, e)$ on the input/output characters tell us the hierarchical position.
The \emph{singlet HMM} generates these labels according to the MixDom stationary process.
The WFST must respect them: it cannot change $(g, e)$ on matched characters.

So the WFST sees:
\begin{itemize}
\item Characters labeled $g=0$: mid-fragment
\item Characters labeled $g=1, e=0$: end of fragment, not end of domain
\item Characters labeled $g=1, e=1$: end of fragment and end of domain
\end{itemize}

At a fragment boundary in the ancestor ($g=1$), the WFST knows the ancestral fragment
has ended. The next ancestral character (if any) will start a new fragment or domain.
Between the end of one fragment and the start of the next, the WFST may insert
new fragments (for $\mathrm{I_F}$) or entire domains (for $\mathrm{I_D}$).

This gives us the complete transition logic from each wait state.

\bigskip
\noindent\textbf{Revised: $\mathrm{W_M}$ at $g=1, e=0$ (fragment boundary, mid-domain):}

\begin{center}
\renewcommand{\arraystretch}{1.4}
\begin{tabular}{llll}
\hline
Dest & Weight & I/O & Notes \\
\hline
$(\mat, \frag', g', \dom, e')$ &
  $\alpha_\dom \cdot \fragdist_{\dom\frag'}$ &
  $\anctok / \destok$ &
  new fragment, match \\
$(\mathrm{D_F}, \frag', g', \dom, e')$ &
  $(1-\alpha_\dom) \cdot \fragdist_{\dom\frag'}$ &
  $\anctok / \varepsilon$ &
  new fragment, delete \\
$(\mathrm{I_F}, \frag', g', \dom, 0)$ &
  $\beta_\dom \cdot \fragdist_{\dom\frag'}$ &
  $\varepsilon / \destok$ &
  insert fragment \\
\hline
\end{tabular}
\end{center}

Normalization: with an input character, the weight is
$\alpha_\dom + (1-\alpha_\dom) = 1$ (times $\fragdist$).
Without an input character (insertion), $\beta_\dom$ (times $\fragdist$).
These don't need to sum to 1 together because input-present and input-absent
are exclusive events in the WFST.

\bigskip
\noindent\textbf{$\mathrm{W_M}$ at $g=1, e=1$ (domain boundary):}

Here both the fragment and domain have ended.
The domain-level TKF91 process decides what happens next.

\begin{center}
\renewcommand{\arraystretch}{1.4}
\begin{tabular}{llll}
\hline
Dest & Weight & I/O & Notes \\
\hline
$(\mat, \frag', g', \dom', e')$ &
  $\nonemptytrans_{\mat\mat} \cdot \domdist_{\dom'} \cdot
   \tkftrans^{(\dom')}_{\sta\mat} \cdot \fragdist_{\dom'\frag'}$ &
  $\anctok / \destok$ &
  new domain, match \\
$(\mathrm{D_F}, \frag', g', \dom', e')$ &
  $\nonemptytrans_{\mat\mat} \cdot \domdist_{\dom'} \cdot
   \tkftrans^{(\dom')}_{\sta\del} \cdot \fragdist_{\dom'\frag'}$ &
  $\anctok / \varepsilon$ &
  new domain, del-frag \\
$(\mathrm{D_D}, \frag', g', \dom', e')$ &
  $\nonemptytrans_{\mat\del} \cdot \domdist_{\dom'} \cdot
   \kappa_{\dom'} \cdot \fragdist_{\dom'\frag'}$ &
  $\anctok / \varepsilon$ &
  delete domain \\
$(\mathrm{I_D}, \frag', g', \dom', e')$ &
  $\nonemptytrans_{\mat\ins} \cdot \domdist_{\dom'} \cdot
   \kappa_{\dom'} \cdot \fragdist_{\dom'\frag'}$ &
  $\varepsilon / \destok$ &
  insert domain \\
$\fin$ &
  $\nonemptytrans_{\mat\fin}$ &
  $\varepsilon / \varepsilon$ &
  end \\
\hline
\end{tabular}
\end{center}

Here $\nonemptytrans_{\mat\cdot}$ are the null-eliminated domain-level transitions.
The $\tkftrans^{(\dom')}_{\sta\mat}$ and $\tkftrans^{(\dom')}_{\sta\del}$ factors
are the TKF92 entry transitions into the new domain $\dom'$.

\bigskip
\noindent\textbf{$\mathrm{W_{D_F}}$ (Wait after Delete-Fragment).}

$\mathrm{W_{D_F}}$ tracks the deletion of individual fragments within a domain
(the domain itself is matched/surviving; only some fragments are deleted).

\medskip
\noindent Case $g=0$ (mid-fragment):
\begin{center}
\renewcommand{\arraystretch}{1.3}
\begin{tabular}{llll}
\hline
Dest & Weight & I/O & Notes \\
\hline
$(\mathrm{D_F}, \frag, g', \dom, e)$ &
  $1$ &
  $\anctok / \varepsilon$ &
  continue deleting fragment \\
\hline
\end{tabular}
\end{center}

\noindent Case $g=1, e=0$ (fragment boundary, mid-domain):
\begin{center}
\renewcommand{\arraystretch}{1.3}
\begin{tabular}{llll}
\hline
Dest & Weight & I/O & Notes \\
\hline
$(\mat, \frag', g', \dom, e')$ &
  $\alpha_\dom \cdot \fragdist_{\dom\frag'}$ &
  $\anctok / \destok$ &
  new frag, match \\
$(\mathrm{D_F}, \frag', g', \dom, e')$ &
  $(1-\alpha_\dom) \cdot \fragdist_{\dom\frag'}$ &
  $\anctok / \varepsilon$ &
  new frag, delete \\
\hline
\end{tabular}
\end{center}

\noindent Case $g=1, e=1$ (domain boundary):
\begin{center}
\renewcommand{\arraystretch}{1.3}
\begin{tabular}{llll}
\hline
Dest & Weight & I/O & Notes \\
\hline
$(\mat, \frag', g', \dom', e')$ &
  $\nonemptytrans_{\mat\mat} \cdot \domdist_{\dom'} \cdot
   \tkftrans^{(\dom')}_{\sta\mat} \cdot \fragdist_{\dom'\frag'}$ &
  $\anctok / \destok$ &
  new domain, match \\
$(\mathrm{D_F}, \frag', g', \dom', e')$ &
  $\nonemptytrans_{\mat\mat} \cdot \domdist_{\dom'} \cdot
   \tkftrans^{(\dom')}_{\sta\del} \cdot \fragdist_{\dom'\frag'}$ &
  $\anctok / \varepsilon$ &
  new domain, del-frag \\
$(\mathrm{D_D}, \frag', g', \dom', e')$ &
  $\nonemptytrans_{\mat\del} \cdot \domdist_{\dom'} \cdot
   \kappa_{\dom'} \cdot \fragdist_{\dom'\frag'}$ &
  $\anctok / \varepsilon$ &
  delete domain \\
$(\mathrm{I_D}, \frag', g', \dom', e')$ &
  $\nonemptytrans_{\mat\ins} \cdot \domdist_{\dom'} \cdot
   \kappa_{\dom'} \cdot \fragdist_{\dom'\frag'}$ &
  $\varepsilon / \destok$ &
  insert domain \\
$\fin$ &
  $\nonemptytrans_{\mat\fin}$ &
  $\varepsilon / \varepsilon$ &
  end \\
\hline
\end{tabular}
\end{center}

Note: $\mathrm{W_{D_F}}$ at domain boundaries uses the \emph{same} $\nonemptytrans_{\mat\cdot}$
row as $\mathrm{W_M}$, because fragment-level deletion within a domain does not affect
the domain-level state (the domain was matched, i.e.\ the top-level state was $\mat$).

\bigskip
\noindent\textbf{$\mathrm{W_{D_D}}$ (Wait after Delete-Domain).}

$\mathrm{W_{D_D}}$ handles the deletion of entire domains.
Within a deleted domain, all fragments are consumed without output.

\medskip
\noindent Case $g=0$ (mid-fragment):
\begin{center}
\renewcommand{\arraystretch}{1.3}
\begin{tabular}{llll}
\hline
Dest & Weight & I/O & Notes \\
\hline
$(\mathrm{D_D}, \frag, g', \dom, e)$ &
  $1$ &
  $\anctok / \varepsilon$ &
  continue consuming domain \\
\hline
\end{tabular}
\end{center}

\noindent Case $g=1, e=0$ (fragment boundary, mid-domain):
\begin{center}
\renewcommand{\arraystretch}{1.3}
\begin{tabular}{llll}
\hline
Dest & Weight & I/O & Notes \\
\hline
$(\mathrm{D_D}, \frag', g', \dom, e')$ &
  $\fragdist_{\dom\frag'}$ &
  $\anctok / \varepsilon$ &
  next fragment in deleted domain \\
\hline
\end{tabular}
\end{center}

Note: within a deleted domain, the entire domain is being consumed,
so the fragment distribution weight $\fragdist_{\dom\frag'}$ is needed to
match the singlet prior.

\noindent Case $g=1, e=1$ (domain boundary):
\begin{center}
\renewcommand{\arraystretch}{1.3}
\begin{tabular}{llll}
\hline
Dest & Weight & I/O & Notes \\
\hline
$(\mat, \frag', g', \dom', e')$ &
  $\nonemptytrans_{\del\mat} \cdot \domdist_{\dom'} \cdot
   \tkftrans^{(\dom')}_{\sta\mat} \cdot \fragdist_{\dom'\frag'}$ &
  $\anctok / \destok$ &
  new domain, match \\
$(\mathrm{D_F}, \frag', g', \dom', e')$ &
  $\nonemptytrans_{\del\mat} \cdot \domdist_{\dom'} \cdot
   \tkftrans^{(\dom')}_{\sta\del} \cdot \fragdist_{\dom'\frag'}$ &
  $\anctok / \varepsilon$ &
  new domain, del-frag \\
$(\mathrm{D_D}, \frag', g', \dom', e')$ &
  $\nonemptytrans_{\del\del} \cdot \domdist_{\dom'} \cdot
   \kappa_{\dom'} \cdot \fragdist_{\dom'\frag'}$ &
  $\anctok / \varepsilon$ &
  delete another domain \\
$(\mathrm{I_D}, \frag', g', \dom', e')$ &
  $\nonemptytrans_{\del\ins} \cdot \domdist_{\dom'} \cdot
   \kappa_{\dom'} \cdot \fragdist_{\dom'\frag'}$ &
  $\varepsilon / \destok$ &
  insert domain \\
$\fin$ &
  $\nonemptytrans_{\del\fin}$ &
  $\varepsilon / \varepsilon$ &
  end \\
\hline
\end{tabular}
\end{center}

Here $\nonemptytrans_{\del\cdot}$ is used because the domain-level state was $\del$.

\bigskip
\noindent\textbf{Insert states ($\mathrm{I_F}$, $\mathrm{I_D}$).}

Insert states handle character-level emissions for inserted fragments/domains.
They self-loop for fragment extension and terminate back to wait states.

\medskip
\noindent $\mathrm{I_F}$ (inserted fragment within an existing domain):
\begin{center}
\renewcommand{\arraystretch}{1.3}
\begin{tabular}{llll}
\hline
Source & Dest & Weight & I/O \\
\hline
$(\mathrm{I_F}, \srcfrag, 0, \dom, e)$ &
  $(\mathrm{I_F}, \destfrag, g', \dom, e)$ &
  $\ext^{(\dom)}_{\srcfrag\destfrag}$ (emit next char) &
  $\varepsilon / \destok$ \\
$(\mathrm{I_F}, \frag, 1, \dom, e)$ &
  $(\mathrm{W_M}, \frag, 1, \dom, e)$ &
  $1$ (fragment ended) &
  $\varepsilon / \varepsilon$ \\
\hline
\end{tabular}
\end{center}

Since the label $g$ on the output character indicates whether the fragment continues ($g=0$)
or ends ($g=1$), the fragment-extension matrix $\ext^{(\dom)}$ is again handled by
the singlet HMM generating the output labels.
In the WFST, $\mathrm{I_F}$ at $g=0$ continues emitting; at $g=1$ the fragment ends
and control returns to $\mathrm{W_M}$ to decide on the next fragment-level action.

Actually, for insert states, the output character labels are generated by the WFST itself
(there is no ancestral character to copy from).
The WFST must therefore assign the correct probabilities to the output labels.
The emission weight at $(\mathrm{I_F}, \frag, g, \dom, e)$ for output character
$\destok_{\class'\frag' g'\dom' e'}$ is:
\begin{equation}
\label{eq:if-emit}
\delta_{\dom\dom'} \delta_{ee'} \cdot
\classdist_{\frag'\class'} \eqm_{\class'\destok} \cdot
\begin{cases}
\ext^{(\dom)}_{\frag\frag'} & \text{if } g' = 0 \\
\notext^{(\dom)}_\frag\,\delta_{\frag\frag'} & \text{if } g' = 1
\end{cases}
\end{equation}
followed by a transition: if $g'=0$, loop to $(\mathrm{I_F}, \frag, 0, \dom, e)$;
if $g'=1$, go to the appropriate wait state.

\medskip
\noindent $\mathrm{I_D}$ (inserted domain): similar to $\mathrm{I_F}$, but the entire
domain is new. Within the inserted domain, the fragment-level TKF92 structure applies:
\begin{center}
\renewcommand{\arraystretch}{1.3}
\begin{tabular}{lll}
\hline
Source condition & Dest & Notes \\
\hline
$(\mathrm{I_D}, \frag, 0, \dom, e)$ &
  $(\mathrm{I_D}, \frag, g', \dom, e)$, emit $\destok$ &
  mid-fragment, continue \\
$(\mathrm{I_D}, \frag, 1, \dom, 0)$ &
  $(\mathrm{I_D}, \frag', g', \dom, e')$, emit $\destok$ &
  new fragment in inserted domain \\
  & weight: $\kappa_\dom \cdot \fragdist_{\dom\frag'}$ & \\
$(\mathrm{I_D}, \frag, 1, \dom, 0)$ &
  domain ends within inserted domain &
  weight: $(1-\kappa_\dom)$ \\
  & $\to$ next domain-level action & \\
$(\mathrm{I_D}, \frag, 1, \dom, 1)$ &
  inserted domain complete &
  returns to domain-level wait \\
\hline
\end{tabular}
\end{center}

When $\mathrm{I_D}$ completes (the inserted domain ends), control goes back to a
domain-level wait state.
Since the domain was inserted (top-level state $\ins$), the next domain-level transition
uses $\nonemptytrans_{\ins\cdot}$:
\begin{center}
\renewcommand{\arraystretch}{1.3}
\begin{tabular}{llll}
\hline
Dest & Weight & I/O & Notes \\
\hline
$(\mat, \frag', g', \dom', e')$ &
  $\nonemptytrans_{\ins\mat} \cdot \domdist_{\dom'} \cdot
   \tkftrans^{(\dom')}_{\sta\mat} \cdot \fragdist_{\dom'\frag'}$ &
  $\anctok / \destok$ &
  match next domain \\
$(\mathrm{D_D}, \frag', g', \dom', e')$ &
  $\nonemptytrans_{\ins\del} \cdot \domdist_{\dom'} \cdot
   \kappa_{\dom'} \cdot \fragdist_{\dom'\frag'}$ &
  $\anctok / \varepsilon$ &
  delete next domain \\
$(\mathrm{I_D}, \frag', g', \dom', e')$ &
  $\nonemptytrans_{\ins\ins} \cdot \domdist_{\dom'} \cdot
   \kappa_{\dom'} \cdot \fragdist_{\dom'\frag'}$ &
  $\varepsilon / \destok$ &
  insert another domain \\
$\fin$ &
  $\nonemptytrans_{\ins\fin}$ &
  $\varepsilon / \varepsilon$ &
  end \\
\hline
\end{tabular}
\end{center}


\subsection{Emission Weights}
\label{sec:wfst-emissions}

The emission weights for each emitting state type are:

\begin{center}
\renewcommand{\arraystretch}{1.4}
\begin{tabular}{lll}
\hline
State & Emission & Weight \\
\hline
$\mat$ & input $\anctok_{\class\frag g\dom e}$, output $\destok_{\class\frag g\dom e}$ &
  $\exp(\revsub_\class \evoltime)_{\anctok\destok}$ \\
$\mathrm{I_F}$ & output $\destok_{\class\frag g\dom e}$ &
  $\classdist_{\frag\class}\, \eqm_{\class\destok} \cdot p_g(\ext^{(\dom)}_{\frag\cdot})$ \\
$\mathrm{I_D}$ & output $\destok_{\class\frag g\dom e}$ &
  $\classdist_{\frag\class}\, \eqm_{\class\destok} \cdot p_g(\ext^{(\dom)}_{\frag\cdot})$ \\
$\mathrm{D_F}$ & input $\anctok_{\class\frag g\dom e}$ &
  $1$ \\
$\mathrm{D_D}$ & input $\anctok_{\class\frag g\dom e}$ &
  $1$ \\
\hline
\end{tabular}
\end{center}
where $p_g(\ext^{(\dom)}_{\frag\cdot})$ stands for $\ext^{(\dom)}_{\frag\frag'}$
if $g=0$ (fragment continues with destination fragchar $\frag'$) and
$\notext^{(\dom)}_\frag$ if $g=1$
(fragment ends), and the substitution matrix $\revsub_\class$ is the rate matrix for site class $\class$.

For the WFST (conditional on ancestor), the delete emission weight is 1 because the
ancestral emission probability is divided out (it was generated by the singlet HMM).
The match emission divides out the ancestral $\eqm_{\class\anctok}$ from the pair emission
$\eqm_{\class\anctok} \exp(\revsub_\class \evoltime)_{\anctok\destok}$,
yielding just the substitution matrix entry.

For insert states, the emission includes the full descendant character probability
because there is no ancestral character to condition on.


\subsection{Verification: Composition Reproduces the Pair HMM}
\label{sec:verification}

\paragraph{Claim.}
The Labeled-MixDom Singlet HMM composed with the Labeled-MixDom WFST
is equivalent to the MixDom Pair HMM defined in Section~\ref{sec:nestpairhmm}.

\paragraph{Sketch of proof.}
The composition proceeds as follows:
\begin{enumerate}
\item The Singlet HMM generates the ancestral labeled sequence, with transitions
  governed by the MixDom stationary distribution.
\item The WFST reads the ancestral labeled sequence as input and produces
  the descendant labeled sequence as output.
\item The composed machine has states that are pairs (singlet state, WFST state).
  Since both are order-1 machines tracking structural labels, the composed state
  is a pair of structural labels.
\end{enumerate}

We verify that the composed transition weights match $\transnest$ for each case:

\paragraph{Match-to-Match within a fragment ($g=0$).}
The singlet emits character $\anctok_{\class\frag 0\dom e}$ with weight
$\classdist_{\frag\class} \eqm_{\class\anctok}$,
and transitions to state $(\frag', g', \dom, e)$ with weight $\ext^{(\dom)}_{\frag\frag'}$ (if $g'=0$)
or $\notext^{(\dom)}_\frag \cdot \ldots$ (if $g'=1$).
The WFST in state $(\mathrm{W_M}, \srcfrag, 0, \dom, e)$ reads this input and
transitions to $(\mat, \destfrag, g', \dom, e)$ with weight $\alpha_\dom$,
emitting $\destok$ with substitution weight $\exp(\revsub_\class \evoltime)_{\anctok\destok}$.
The composed emission weight is
$\classdist_{\destfrag\class} \eqm_{\class\anctok} \exp(\revsub_\class \evoltime)_{\anctok\destok}$,
summing over $\class$ gives
$\sum_\class \classdist_{\destfrag\class} \eqm_{\class\anctok} \exp(\revsub_\class \evoltime)_{\anctok\destok}$,
which matches the Pair HMM emission for $\mat\mat_{\dom\destfrag}$.
The transition weight within the fragment is $\ext^{(\dom)}_{\srcfrag\destfrag}
\cdot \alpha_\dom$, which corresponds to the $\samedom \ext^{(\dom)}_{\srcfrag\destfrag}$
term in $\transnest$ (for the intra-fragment Markov contribution to
$\mat\mat_{\srcdom\srcfrag} \to \mat\mat_{\srcdom\destfrag}$).

\paragraph{Match-to-Match across fragment boundary ($g=1, e=0$).}
Singlet: fragment ends, transitions to new fragment $\frag'$ within same domain
with weight $\kappa_\dom \cdot \fragdist_{\dom\frag'}$.
WFST: $(\mathrm{W_M}, \frag, 1, \dom, 0)$ transitions to $(\mat, \frag', g', \dom, e')$
with weight $\alpha_\dom \cdot \fragdist_{\dom\frag'}$.
Combined: $\kappa_\dom \cdot \fragdist_{\dom\frag'} \cdot \alpha_\dom \cdot \fragdist_{\dom\frag'}$.

Wait---the $\fragdist$ appears twice, which is wrong.
This reveals that the WFST transition weight should \emph{not} include $\fragdist_{\dom\frag'}$
when the input character determines $\frag'$.
The fragment type of the next input character is determined by the singlet HMM,
and the WFST simply reads whatever fragment type appears.

Let us correct: at fragment boundaries, when the next action involves reading an input
character, the WFST does \emph{not} weight by $\fragdist_{\dom\frag'}$.
The $\fragdist$ is part of the singlet distribution, not the conditional (WFST) distribution.

\paragraph{Corrected Wait-State Transitions.}
The WFST represents the \emph{conditional} distribution $P(\text{descendant} \mid \text{ancestor})$.
Therefore:
\begin{itemize}
\item Transitions that consume an input character should \emph{not} include the prior
  probability of that input character's labels (such as $\fragdist$, $\domdist$).
\item Transitions that produce an output character (insertions) \emph{do} include
  the full probability of the output labels (since the WFST generates them).
\item Transitions involving the $\nonemptytrans$ matrix at domain boundaries must be
  adjusted: $\nonemptytrans$ was derived for the Pair HMM (joint distribution),
  so the WFST version divides out the ancestral prior factors.
\end{itemize}

Concretely, the $\transnest$ entry for
$\mat\xstate_{\srcdom\srcfrag} \to \mat\ystate_{\destdom\destfrag}$ with $\srcdom = \destdom$
(same domain, new fragment) is:
\[
\notext^{(\dom)}_\srcfrag \cdot \tkftrans^{(\dom)}_{\xstate\ystate} \cdot \fragdist_{\dom\destfrag}
\]
This is the \emph{joint} weight.
In the composed (singlet $\circ$ WFST) machine:
\begin{itemize}
\item Singlet provides: $\notext^{(\dom)}_\srcfrag \cdot \kappa_\dom \cdot \fragdist_{\dom\destfrag}$
  (end fragment, continue domain, choose new fragment type)
\item WFST provides: $\frac{\tkftrans^{(\dom)}_{\xstate\ystate}}{\kappa_\dom}$
  (the conditional transition, dividing out the $\kappa_\dom$ from the joint)
\end{itemize}
Product: $\notext^{(\dom)}_\srcfrag \cdot \kappa_\dom \cdot \fragdist_{\dom\destfrag}
\cdot \tkftrans^{(\dom)}_{\xstate\ystate} / \kappa_\dom
= \notext^{(\dom)}_\srcfrag \cdot \tkftrans^{(\dom)}_{\xstate\ystate} \cdot \fragdist_{\dom\destfrag}$.
This matches the Pair HMM. \checkmark

Similarly, for the domain-boundary case with $\srcdom \neq \destdom$:
\begin{itemize}
\item Singlet provides: $\notext^{(\srcdom)}_\srcfrag (1-\kappa_{\srcdom}) \cdot
  \frac{\kappa_\main \domdist_{\destdom} \kappa_{\destdom} \fragdist_{\destdom\destfrag}}{1-\zeta}$
\item WFST provides: the conditional weight that, when multiplied by the singlet weight,
  gives $\transnest_{\mat\xstate_{\srcdom\srcfrag} \to \mat\ystate_{\destdom\destfrag}}$
\end{itemize}

The required WFST transition weights therefore depend on the singlet transition structure.
Rather than writing out all the corrected weights with singlet factors divided out,
we express the key principle:

\begin{quote}
\textbf{WFST weight principle:}
For any transition that consumes an input character with label $\ell'$,
the WFST weight is the Pair HMM transition weight divided by the singlet
transition weight for generating a character with label $\ell'$ from the current
singlet state.
For transitions that produce an output character (insertions), the WFST
carries the full conditional weight.
For null transitions (emitting-to-wait), the weight is 1.
\end{quote}

Specifically, let $P_{\rm pair}(i \to j)$ be the Pair HMM transition from state $i$ to $j$,
and let $P_{\rm sing}(\ell \to \ell')$ be the singlet transition.
Then:
\begin{equation}
\label{eq:wfst-weight-principle}
w_{\rm WFST}(s, \ell_{\rm in}, \ell_{\rm out} \to s', \ell') =
\frac{P_{\rm pair}(i(s,\ell) \to j(s',\ell'))}{P_{\rm sing}(\ell \to \ell')}
\end{equation}
where $i(s,\ell)$ and $j(s',\ell')$ are the corresponding Pair HMM states,
and $\ell_{\rm in}$ is present when $s'$ consumes input.

\paragraph{Explicit corrected weights.}

We now tabulate the corrected WFST transition weights organized by source wait state,
boundary case, and whether input is consumed.

In all cases below, the WFST context $\ell = (\frag, g, \dom, e)$ is the structural label
of the last processed character. The singlet HMM is in the corresponding state.
The factor $P_{\rm sing}(\ell \to \ell')$ is the singlet transition weight from the
current label to the next label.

\bigskip
\noindent\textbf{$\mathrm{W_M}$, $g=0$ (mid-fragment, input consumed):}

\begin{center}
\renewcommand{\arraystretch}{1.4}
\begin{tabular}{lll}
\hline
Dest & WFST weight & Notes \\
\hline
$\mat$ & $\alpha_\dom$ & match \\
$\mathrm{D_F}$ & $(1-\alpha_\dom)$ & frag-level delete \\
\hline
\end{tabular}
\end{center}
No singlet factor to divide out: at $g=0$ the singlet continues the fragment
with weight $\ext^{(\dom)}_{\srcfrag\destfrag}$ (producing the next character
with destination fragchar $\destfrag$ in the same domain), and the WFST
weight is purely the TKF92 match/delete split, independent of $\fragdist$.

\bigskip
\noindent\textbf{$\mathrm{W_M}$, $g=1, e=0$ (fragment boundary, mid-domain):}

With input (start new ancestral fragment):
\begin{center}
\renewcommand{\arraystretch}{1.4}
\begin{tabular}{lll}
\hline
Dest & WFST weight & Notes \\
\hline
$\mat$ & $\alpha_\dom$ & match new fragment \\
$\mathrm{D_F}$ & $(1-\alpha_\dom)$ & delete new fragment \\
\hline
\end{tabular}
\end{center}
The $\fragdist_{\dom\frag'}$ is supplied by the singlet; the WFST contributes
only the match/delete split.

Without input (descendant insertion):
\begin{center}
\renewcommand{\arraystretch}{1.4}
\begin{tabular}{lll}
\hline
Dest & WFST weight & Notes \\
\hline
$(\mathrm{I_F}, \frag', g', \dom, 0)$ &
  $\beta_\dom \cdot \fragdist_{\dom\frag'} \cdot \classdist_{\frag'\class'} \eqm_{\class'\destok} \cdot p_{g'}(\ext_{\frag'})$ &
  insert fragment \\
\hline
\end{tabular}
\end{center}
Here the WFST generates the full output character probability since
there is no input to condition on.

\bigskip
\noindent\textbf{$\mathrm{W_M}$, $g=1, e=1$ (domain boundary):}

With input (start new ancestral domain):
\begin{center}
\renewcommand{\arraystretch}{1.4}
\begin{tabular}{lll}
\hline
Dest & WFST weight & Notes \\
\hline
$\mat$ &
  $\dfrac{\nonemptytrans_{\mat\mat} \cdot \domdist_{\dom'} \cdot
   \tkftrans^{(\dom')}_{\sta\ystate} \cdot \fragdist_{\dom'\frag'}}
  {P_{\rm sing}(\ell \to \ell')}$ &
  new domain, match \\[8pt]
$\mathrm{D_F}$ &
  $\dfrac{\nonemptytrans_{\mat\mat} \cdot \domdist_{\dom'} \cdot
   \tkftrans^{(\dom')}_{\sta\del} \cdot \fragdist_{\dom'\frag'}}
  {P_{\rm sing}(\ell \to \ell')}$ &
  new domain, del-frag \\[8pt]
$\mathrm{D_D}$ &
  $\dfrac{\nonemptytrans_{\mat\del} \cdot \domdist_{\dom'} \cdot
   \kappa_{\dom'} \cdot \fragdist_{\dom'\frag'}}
  {P_{\rm sing}(\ell \to \ell')}$ &
  delete domain \\
\hline
\end{tabular}
\end{center}

Without input:
\begin{center}
\renewcommand{\arraystretch}{1.4}
\begin{tabular}{lll}
\hline
Dest & WFST weight & Notes \\
\hline
$\mathrm{I_D}$ &
  $\nonemptytrans_{\mat\ins} \cdot \domdist_{\dom'} \cdot
   \kappa_{\dom'} \cdot \fragdist_{\dom'\frag'} \cdot
   \classdist_{\frag'\class'} \eqm_{\class'\destok} \cdot p_{g'}(\ext_{\frag'})$ &
  insert domain \\
$\fin$ & $\nonemptytrans_{\mat\fin} / P_{\rm sing}(\ell \to \fin)$ &
  end \\
\hline
\end{tabular}
\end{center}

where $P_{\rm sing}(\ell \to \ell') = \frac{\kappa_\main \domdist_{\dom'} \kappa_{\dom'} \fragdist_{\dom'\frag'}}{1-\zeta}$
is the singlet transition from a domain boundary to the next character label $\ell'$
(Equation~\ref{eq:singlet-newdom}).

\bigskip
The $\mathrm{W_{D_F}}$ and $\mathrm{W_{D_D}}$ tables follow the same pattern,
using $\nonemptytrans_{\mat\cdot}$ for $\mathrm{W_{D_F}}$ (since fragment deletion
is within a matched domain) and $\nonemptytrans_{\del\cdot}$ for $\mathrm{W_{D_D}}$
(since the domain is being deleted).

\paragraph{$\mathrm{W_{D_F}}$ at $g=0$:}
\begin{center}
\renewcommand{\arraystretch}{1.3}
\begin{tabular}{ll}
\hline
Dest & WFST weight \\
\hline
$\mathrm{D_F}$ & $1$ \\
\hline
\end{tabular}
\end{center}

\paragraph{$\mathrm{W_{D_F}}$ at $g=1, e=0$:}
\begin{center}
\renewcommand{\arraystretch}{1.3}
\begin{tabular}{ll}
\hline
Dest & WFST weight \\
\hline
$\mat$ & $\alpha_\dom$ \\
$\mathrm{D_F}$ & $(1-\alpha_\dom)$ \\
\hline
\end{tabular}
\end{center}

\paragraph{$\mathrm{W_{D_F}}$ at $g=1, e=1$:}
Same structure as $\mathrm{W_M}$ at $g=1, e=1$, using $\nonemptytrans_{\mat\cdot}$.

\paragraph{$\mathrm{W_{D_D}}$ at $g=0$:}
\begin{center}
\renewcommand{\arraystretch}{1.3}
\begin{tabular}{ll}
\hline
Dest & WFST weight \\
\hline
$\mathrm{D_D}$ & $1$ \\
\hline
\end{tabular}
\end{center}

\paragraph{$\mathrm{W_{D_D}}$ at $g=1, e=0$:}
\begin{center}
\renewcommand{\arraystretch}{1.3}
\begin{tabular}{ll}
\hline
Dest & WFST weight \\
\hline
$\mathrm{D_D}$ & $1$ \\
\hline
\end{tabular}
\end{center}
(Within a deleted domain, all fragments are consumed; fragment type is irrelevant.)

\paragraph{$\mathrm{W_{D_D}}$ at $g=1, e=1$:}
\begin{center}
\renewcommand{\arraystretch}{1.3}
\begin{tabular}{lll}
\hline
Dest & WFST weight & Notes \\
\hline
$\mat$ &
  $\nonemptytrans_{\del\mat} \cdot \domdist_{\dom'} \cdot
   \tkftrans^{(\dom')}_{\sta\ystate} \cdot \fragdist_{\dom'\frag'}
   / P_{\rm sing}(\ell \to \ell')$ &
  new domain, match \\
$\mathrm{D_F}$ &
  $\nonemptytrans_{\del\mat} \cdot \domdist_{\dom'} \cdot
   \tkftrans^{(\dom')}_{\sta\del} \cdot \fragdist_{\dom'\frag'}
   / P_{\rm sing}(\ell \to \ell')$ &
  new domain, del-frag \\
$\mathrm{D_D}$ &
  $\nonemptytrans_{\del\del} \cdot \domdist_{\dom'} \cdot
   \kappa_{\dom'} \cdot \fragdist_{\dom'\frag'}
   / P_{\rm sing}(\ell \to \ell')$ &
  delete domain \\
$\mathrm{I_D}$ &
  $\nonemptytrans_{\del\ins} \cdot (\text{full output prob})$ &
  insert domain \\
$\fin$ & $\nonemptytrans_{\del\fin} / P_{\rm sing}(\ell \to \fin)$ &
  end \\
\hline
\end{tabular}
\end{center}

\paragraph{$\mathrm{I_D}$ completion at $g=1, e=1$:}
\begin{center}
\renewcommand{\arraystretch}{1.3}
\begin{tabular}{lll}
\hline
Dest & WFST weight & Notes \\
\hline
$\mat$ &
  $\nonemptytrans_{\ins\mat} \cdot \domdist_{\dom'} \cdot
   \tkftrans^{(\dom')}_{\sta\ystate} \cdot \fragdist_{\dom'\frag'}
   / P_{\rm sing}(\ell \to \ell')$ &
  match next domain \\
$\mathrm{D_D}$ &
  $\nonemptytrans_{\ins\del} \cdot \domdist_{\dom'} \cdot
   \kappa_{\dom'} \cdot \fragdist_{\dom'\frag'}
   / P_{\rm sing}(\ell \to \ell')$ &
  delete next domain \\
$\mathrm{I_D}$ &
  $\nonemptytrans_{\ins\ins} \cdot (\text{full output prob})$ &
  insert another domain \\
$\fin$ & $\nonemptytrans_{\ins\fin} / P_{\rm sing}(\ell \to \fin)$ &
  end \\
\hline
\end{tabular}
\end{center}


\subsection{Simplification of Domain-Boundary WFST Weights}
\label{sec:wfst-simplify}
\label{sec:simplification}

The domain-boundary WFST weights (Section~\ref{sec:wait-outgoing}) involve
a ratio of the Pair HMM transition weight to the singlet transition weight.
This simplifies considerably.

For a transition that consumes input with label $\ell' = (\frag', g', \dom', e')$
from a domain boundary ($g=1, e=1$), the singlet weight is
(from Equation~\ref{eq:singlet-newdom}):
\[
P_{\rm sing}(\ell \to \ell') = \frac{\kappa_\main \cdot \domdist_{\dom'} \cdot \kappa_{\dom'} \cdot \fragdist_{\dom'\frag'}}{1 - \zeta}
\]

The Pair HMM weight for $\mat\mat_{\srcdom\srcfrag} \to \mat\mat_{\destdom\destfrag}$
with $\srcdom \neq \destdom$ (different domain, going through domain boundary) is:
\[
\notext^{(\srcdom)}_\srcfrag \cdot \tkftrans^{(\srcdom)}_{\mat\fin} \cdot
\nonemptytrans_{\mat\mat} \cdot \domdist_{\destdom} \cdot
\tkftrans^{(\destdom)}_{\sta\mat} \cdot \fragdist_{\destdom\destfrag}
\]

The WFST weight is therefore:
\[
\frac{\notext^{(\srcdom)}_\srcfrag \cdot \tkftrans^{(\srcdom)}_{\mat\fin} \cdot
\nonemptytrans_{\mat\mat} \cdot \domdist_{\destdom} \cdot
\tkftrans^{(\destdom)}_{\sta\mat} \cdot \fragdist_{\destdom\destfrag}}
{\frac{\kappa_\main \cdot \domdist_{\destdom} \cdot \kappa_{\destdom} \cdot \fragdist_{\destdom\destfrag}}{1 - \zeta}}
= \frac{(1-\zeta)\notext^{(\srcdom)}_\srcfrag \cdot \tkftrans^{(\srcdom)}_{\mat\fin} \cdot
\nonemptytrans_{\mat\mat} \cdot \tkftrans^{(\destdom)}_{\sta\mat}}
{\kappa_\main \cdot \kappa_{\destdom}}
\]

The $\domdist_{\destdom}$ and $\fragdist_{\destdom\destfrag}$ cancel, leaving a weight
that depends on the source domain parameters and the destination domain's TKF parameters,
but \emph{not} on the specific fragment or domain type of the destination.
This is a significant simplification: the WFST transition weight at domain boundaries
is the same for all destination labels $\ell'$, given the source label $\ell$.

Since $\tkftrans^{(\srcdom)}_{\mat\fin} = (1-\beta_{\srcdom})(1-\kappa_{\srcdom})$
and $\tkftrans^{(\destdom)}_{\sta\mat} = (1-\beta_{\destdom})\kappa_{\destdom}\alpha_{\destdom}$,
the weight becomes:
\[
\frac{(1-\zeta)\notext^{(\srcdom)}_\srcfrag(1-\beta_{\srcdom})(1-\kappa_{\srcdom}) \cdot
\nonemptytrans_{\mat\mat} \cdot (1-\beta_{\destdom})\alpha_{\destdom}}
{\kappa_\main}
\]

Note that this weight still depends on $\destdom$ through $(1-\beta_{\destdom})\alpha_{\destdom}$,
so the cancellation is partial: $\domdist$ and $\fragdist$ cancel but the TKF92 entry
parameters do not.


\subsection{State Count and Sparsity}
\label{sec:state-count}

The Labeled-MixDom WFST has at most $8L + 2$ states where $L = |\nfrag| \cdot |\ndom| \cdot 4$.
In practice, many combinations are constrained away:
\begin{itemize}
\item $\mathrm{I_F}$ states only occur with $e \neq 1$ (inserted fragments cannot be the last in a domain
  since they are inserted \emph{within} a domain).
\item Mid-fragment ($g=0$) wait states have at most 2 outgoing transitions each
  (continue in same fragment, match or delete).
\item Fragment-boundary ($g=1, e=0$) wait states have at most 3 outgoing transition types.
\item Domain-boundary ($g=1, e=1$) wait states have at most 5 outgoing transition types
  (one per $\nonemptytrans$ column).
\end{itemize}

The transition matrix is therefore very sparse.
For typical values ($|\nfrag| = 4$, $|\ndom| = 20$, $|\nclass| = 4$),
we get $L = 4 \cdot 20 \cdot 4 = 320$ and the WFST has at most $8 \cdot 320 + 2 = 2562$ states,
comparable to the Maraschino-distilled WFST
(which has $\sim 8 \cdot |\alphabet|^2 + 2$ states for the order-1 pair context).

The key advantage over distillation is exactness: the Labeled-MixDom WFST preserves
all correlations of the MixDom model without approximation, at the cost of a larger
effective alphabet.
